\section{Deep Learning Literature}
\label{sec:literature}

The C-MAPSS benchmark was introduced as a run-to-failure simulation of
turbofan degradation \cite{saxena_2008_damage_prognosis}.  Its engine-indexed
multivariate trajectories make it a natural test bed for models that learn
temporal context rather than treating observations as independent rows.
Long short-term memory (LSTM) networks use gated memory to preserve useful
information over long intervals \cite{hochreiter_1997_lstm}; gated recurrent
units (GRUs) provide a smaller alternative \cite{cho_2014_gru}.  In contrast,
temporal convolutional networks (TCNs) use causal convolutions and a finite
receptive field \cite{bai_2018_tcn}, while the Transformer replaces recurrence
with attention \cite{vaswani_2017_transformer}.

RUL studies commonly use convolutional encoders because local changes in
sensor trajectories are predictive.  Babu et al. proposed a deep CNN
regression approach for C-MAPSS \cite{babu_2016_dcnn}, and Li et al. studied
deep convolutional features for prognostics \cite{li_2018_dcnn}.  LSTM-based
RUL work further demonstrated that a sequence representation can model
degradation history \cite{zheng_2017_lstm}.  Attention and CNN--recurrent
hybrids are attractive extensions, but their extra parameters are a material
risk when only 100 engines are available.  The literature also frequently
uses a piecewise-linear cap near 125 cycles: this reduces the influence of
uncertain early-life labels, but can remove information needed to correct the
high-RUL underprediction seen in our uncapped Phase 1 baseline.

\begin{table}[t]
\caption{Scope of the reviewed literature. ``n/a'' means that the source or
the accessible report did not state the item unambiguously.}
\label{tab:literature}
\centering
\scriptsize
\resizebox{\columnwidth}{!}{%
\begin{tabular}{p{0.25\columnwidth}p{0.18\columnwidth}p{0.15\columnwidth}p{0.28\columnwidth}}
\hline
Source & Model & Subset & Limitation \\
\hline
Babu et al. & CNN & FD001--FD004 & Representation-sensitive \\
Li et al. & CNN & FD001--FD004 & Preprocessing-sensitive \\
Zheng et al. & LSTM & FD001--FD004 & Capping discards high-RUL detail \\
Bai et al. & TCN & n/a & Receptive-field dependent \\
Vaswani et al. & Transformer & n/a & Attention cost grows with length \\
\hline
\end{tabular}}
\end{table}

Unlike Phase 1's ten-cycle rolling summaries, a window model receives the
ordered sensor values and can learn interactions between trends, transients,
and delayed effects.  Reported windows vary by task, and there is no
universally optimal length: short windows miss slow degradation while long
windows increase padding, computation, and overfitting.  These observations
motivate a controlled comparison of recurrent, convolutional, attention, and
hybrid representations on the same engine split.
